\FloatBarrier
\section{Adapter and torus exit shield}

The remaining downstream contacts recorded in this historical study are resolved by the subsequent socket
revision in Section~\ref{sec:downstream-joints}.

The third beamline study covers the adapter and torus exit-shield casing and lead fill. The reducer is
renamed Beamline\_Adapter\_1 in GEMC3 and GEMC2, including all eleven source CAD variations, generated
runtime GXML, and eleven run-dependent SQLite CAD rows. Every other database field, material, coordinate
frame, and placement is retained; custom GXML must use the new adapter name to match the STL filename.

These three meshes are already coarse, closed, manifold, and free of detected self-intersections and
zero-area triangles. The adapter decreases from 528 to 526 triangles and the casing from 1,392 to 1,384. The
lead fill remains byte-identical at 384 triangles, with both closed components retained: even the smallest
trial collapse displaced its surface by more than 1 mm near tight mating clearances. Its exact geometry was
explicitly retained after reviewing this constraint. The group decreases from 2,304 to 2,294 triangles
(0.434\%) and from 115,452 to 114,952 bytes.

Table~\ref{tab:exit-meshes} summarizes the mesh inventory.

\begin{table}[htbp]
\centering
\caption{Adapter and exit-shield triangle counts.}
\label{tab:exit-meshes}
\begin{tabular}{lrr}
\toprule
Mesh & Original & Optimized \\
\midrule
Adapter & 528 & 526 \\
Exit casing & 1,392 & 1,384 \\
Exit lead fill & 384 & 384 \\
\bottomrule
\end{tabular}
\end{table}

For the casing, 432 faces having at least one vertex within 1 mm of the lead fill are protected. The saved
output is checked to retain every protected triangle at exactly the original coordinates. All three meshes
retain their bounds, component counts, and genus; no hole filling or reconstruction is used. The adapter and
casing have maximum bidirectionally sampled surface differences of 0.4559 and 0.5002 mm, below the 0.51 mm
sampling guard for this group, and volume changes of +0.000368\% and +0.000188\%. The guard is specific to
this coarse group and does not establish the apex study's 0.05 mm criterion. Every vertex and approximately
200,000 edge/surface samples per direction are checked; the sampled maxima are not proven global bounds. The
builder reproduces all three outputs byte for byte from the archived originals.

The benchmark isolates these three volumes with their unchanged aluminum, stainless-steel, and lead
materials in air. The recorded GEMC3 executable and Geant4 stack run FTFP\_BERT without optical processes,
zero magnetic field, one thread, default production cuts, and no GUI or event output. Matched 5 GeV LUND
primaries start at the target, uniform in solid angle over 0.25--2.6 degrees and full azimuth. Five
alternating-order repetitions use seeds 12345--12349 and 200 events per electron, proton, positive-pion, or
neutron run. Separate zero-event runs measure startup without overlap checking; time since beamOn measures
transport. CPU time and peak RSS are retained separately. Absolute timings are not pooled with either
earlier beamline study or LTCC.

Table~\ref{tab:exit-timing} gives the separate transport measurements.

\begin{table}[htbp]
\centering
\caption{Median transport wall seconds per 200 events over five repetitions.}
\label{tab:exit-timing}
\begin{tabular}{lrrr}
\toprule
Particle & Original & Optimized & Original/optimized \\
\midrule
e- & 10.750 & 10.443 & 1.029 \\
proton & 3.451 & 3.398 & 1.016 \\
pi+ & 4.363 & 4.269 & 1.022 \\
neutron & 3.421 & 3.293 & 1.039 \\
\bottomrule
\end{tabular}
\end{table}

Median startup simulation wall time changes from 0.8844 to 0.8707 s, and peak loading RSS from 125.23 to
125.16 MiB. The original/optimized median transport wall-time ratios range from 1.016 to 1.039; individual
seed-paired ratios span 1.015--1.052. With only ten removed triangles and five repetitions, these timings do
not establish a substantial performance gain. All 52 processes exit normally, and all forty transport runs
report the 200 requested events.

Table~\ref{tab:exit-warnings} gives the sampled diagnostic counts.

\begin{table}[htbp]
\centering
\caption{Geant4 exception counts over all 52 isolated group processes.}
\label{tab:exit-warnings}
\begin{tabular}{lrr}
\toprule
Diagnostic & Original & Optimized \\
\midrule
GeomSolids1001 & 0 & 0 \\
GeomNav1002 & 0 & 0 \\
GeomNav0003 & 0 & 0 \\
GeomVol1002 & 0 & 0 \\
\bottomrule
\end{tabular}
\end{table}

No Geant4 exceptions are recorded in the isolated loading, transport, or overlap samples for either case.
This is a limited sample, not proof of complete overlap freedom or detector-response equivalence. Full
default beamline construction also completes with the renamed adapter and no missing mesh; its seven solid
warnings and twenty overlap reports are unchanged from the preceding downstream state. All fifteen native
geometry/material variations and fourteen CAD variations pass the semantic comparison against GEMC2.

\FloatBarrier
Figure~\ref{fig:exit-stl} compares the stored triangulations.

\begin{figure}[htbp]
\centering
\includegraphics[width=0.84\textwidth]{figures/exit_stl_comparison.pdf}
\caption{Original (left) and final (right) adapter and torus exit-shield STL surfaces in millimetres.
Each pair shares camera and bounds; the lead fill is retained exactly.}
\label{fig:exit-stl}
\end{figure}
